\clearpage
\VerdierPhysicalPageAnchor{0074}%
\VerdierSemanticAnchor{ch1-2-4-1-pi-opposite}
(2.2.3) と (2.2.4) で定義された関手の反対関手である。同様に、次のように表す：

\VerdierSemanticAnchor{ch1-2-4-1-eq4}
\begin{equation}
  \int_{\varphi,C^\circ}^{\Pi}\colon
  \operatorname{comp}^{n}(C^\circ)
  \longrightarrow
  \operatorname{comp}^{m}(C^\circ),
  \tag{2.4.1.4}
\end{equation}

\VerdierSemanticAnchor{ch1-2-4-1-pi-definition}
これは (2.2.3) と (2.2.4) で定義された関手である。次の図式が可換であることは
ただちに確かめられる：

\VerdierSemanticAnchor{ch1-2-4-1-eq5}
\begin{equation}
  \resizebox{0.82\textwidth}{!}{%
    $\begin{tikzcd}[ampersand replacement=\&,column sep=6.5em,row sep=5.2em]
      \operatorname{comp}^{n}(C)^\circ
        \arrow[r,"\iota_n"]
        \arrow[d,"{\left(\int_{\varphi,C}^{\Sigma}\right)^\circ}"']
        \& \operatorname{comp}^{n}(C^\circ)
          \arrow[d,"{\int_{\varphi,C^\circ}^{\Pi}}"] \\
      \operatorname{comp}^{m}(C)^\circ
        \arrow[r,"\iota_m"]
        \& \operatorname{comp}^{m}(C^\circ)
    \end{tikzcd}$%
  }
  \tag{2.4.1.5}
\end{equation}

\VerdierSemanticAnchor{ch1-2-4-1-eq5-commutative}
特に、$\varphi$ が唯一の写像
$[n]\longrightarrow[1]$ であるとき、次の可換図式を得る：

\VerdierSemanticAnchor{ch1-2-4-1-eq6}
\begin{equation}
  \resizebox{0.82\textwidth}{!}{%
    $\begin{tikzcd}[ampersand replacement=\&,column sep=6.5em,row sep=5.2em]
      \operatorname{comp}^{n}(C)^\circ
        \arrow[r,"\iota_n"]
        \arrow[d,"{\left(\int_{n,C}^{\Sigma}\right)^\circ}"']
        \& \operatorname{comp}^{n}(C^\circ)
          \arrow[d,"{\int_{n,C^\circ}^{\Pi}}"] \\
      \operatorname{comp}(C)^\circ
        \arrow[r,"\iota_1"]
        \& \operatorname{comp}(C^\circ)
    \end{tikzcd}$%
  }
  \tag{2.4.1.6}
\end{equation}

\VerdierSemanticAnchor{ch1-2-4-1-twisted-isomorphism}
また、同型の組 $(\iota_1,\iota_n)$ が、捩れた
$\mathbb Z[n]$-関手 $\left(\int_{n,C}^{\Sigma}\right)^\circ$ から捩れた
$\mathbb Z[n]$-関手 $\int_{n,C^\circ}^{\Pi}$ への同型であることもただちに確かめられる
（$\int_{n,C^\circ}^{\Pi}$ の次数付き構造は (2.3.4) で述べられ、
$\left(\int_{n,C}^{\Sigma}\right)^\circ$ の次数付き構造は (2.3.4) と
(1.4.8) で述べられる）。

\VerdierSemanticAnchor{ch1-2-4-1-identification}
以下、同型 $\iota_n$ により、圏
$\operatorname{comp}^{n}(C)^\circ$ と
$\operatorname{comp}^{n}(C^\circ)$、および関手
$\left(\int_{\varphi,C}^{\Sigma}\right)^\circ$ と
$\int_{\varphi,C^\circ}^{\Pi}$ をそれぞれ同一視する。

\VerdierSemanticAnchor{ch1-2-4-2}
\noindent\textbf{2.4.2.}\quad
$C_i$（$i\in[n]$）を加法圏の族とし、
$\prod_{i\in[n]}C_i$ をその積とする。
$\prod_{i\in[n]}C_i$ の対象を $(X_1,X_2,\ldots,X_n)$ と表し、記号
$(X_1,X_2,\ldots,f_i,\ldots,X_n)$ で次の射を表す：

\VerdierSemanticAnchor{ch1-2-4-2-component-morphism}
\begin{equation*}
  (\operatorname{id}_{X_1},\operatorname{id}_{X_2},\ldots,
  f_i,\ldots,\operatorname{id}_{X_n}).
\end{equation*}
